\documentclass[11pt]{article}
\usepackage[margin=0.75in]{geometry}
\usepackage{fontspec}
\usepackage{setspace}
\usepackage{hyperref}
\hypersetup{colorlinks=true, urlcolor=blue, linkcolor=black}
\pagestyle{empty}
\setlength{\parskip}{0.55em}
\setlength{\parindent}{0pt}
\setstretch{0.97}
\setmainfont{Times New Roman}

\begin{document}

\begin{flushleft}
Ke Cheng HSU (Clio) \\
Department of Economics, The Hong Kong University of Science and Technology \\
Clear Water Bay, Kowloon, Hong Kong $\cdot$ +852 6840 9753 \\
\href{mailto:khsu@connect.ust.hk}{khsu@connect.ust.hk} $\cdot$ \href{https://cliohsu1997.github.io/}{cliohsu1997.github.io}
\end{flushleft}

2 November 2026

\vspace{0.25em}
Faculty Search Committee \\
Economics and Management School and Institute for Advanced Study, Wuhan University \\
Wuhan, Hubei, China

\vspace{0.25em}
Dear Members of the Search Committee,

I am writing to apply for the tenure-track Assistant Professor position in the Economics and Management School (EMS) and the Institute for Advanced Study (IAS) at Wuhan University (JOE 2026-02\_111478087). I am a Ph.D.\ candidate in Economics at the Hong Kong University of Science and Technology. I expect to graduate in June 2027 and can join in Fall 2027. My primary fields are Empirical Industrial Organization and Digital Economics.

My job market paper, ``Designing Return Policy in E-Commerce: A Trade-off between Purchase and Return Efficiency,'' studies how return policy shapes both the purchase margin and the return margin. A stricter policy cuts wasteful returns after delivery, while lower return rates make a higher product price more profitable for the seller and can price out efficient purchases. Using transaction data from a Taobao apparel retailer across six product categories from 2017 to 2023, I estimate a structural model of consumers' purchase, return-shipping insurance, and return choices under monopoly pricing. A lenient return policy performs better when margins and salvage values are high by expanding valuable purchases, and a strict policy performs better when margins are thinner by curbing excessive returns. Pairing a discounted price with a restocking fee balances purchase and return efficiency and outperforms existing return-shipping insurance for most categories.

In my other projects, with Pak Hung Au, I study how firms use product design to direct the search behavior of consumers with different search costs. With Kohei Kawaguchi and Kei Kawai, I use Japanese bookstore point-of-sale data to study how price commitment encourages retailers to hold inventory while leaving unsold copies that are returned when demand is low. With Ce Bian and Liang Zhao, I study how product returns change tax pass-through and who bears the tax, and show that a tax can lower the equilibrium price and raise welfare.

At HKUST I have taught as a teaching assistant from 2020 to 2025 for undergraduate macroeconomics, the economics of uncertainty and information, and Ph.D.\ microeconomic theory, with an average evaluation of 4.72/5. In 2026 I designed and lectured the mathematics camp for the MSc in Economics, receiving 4.53/5. I am prepared to teach microeconomic theory, industrial organization, applied microeconomics, information economics, and econometrics from the undergraduate to the Ph.D.\ level, mathematics for economists through the master's level, and macroeconomics at the undergraduate level.

I submit my CV, job market paper, and research statement via JOE, and Pak Hung Au, Kohei Kawaguchi, and Kei Kawai will submit reference letters through JOE. You can reach me at khsu@connect.ust.hk or +852 6840 9753.

\vspace{0.6em}
Sincerely,

\vspace{0.6em}
Ke Cheng HSU \\
Ph.D.\ Candidate, Department of Economics, HKUST

\end{document}
